\section{Conclusion}
\label{sec:conclusion}

A11yPreserve provides a source-grounded way to test whether selected accessibility semantics survive DOCX-to-PDF conversion. The frozen V4 design combines feature contracts, raw OOXML verification, a separately implemented source oracle, real conversion pipelines, multiple destination evidence paths, conservative classifications, prospective holdout evaluation, calibration, mutation testing, repeatability, and a separate realistic-document arm.

The prospectively frozen primary evaluation contains 130 LibreOffice and Google Docs fixture--pipeline cases. LibreOffice produced 38 verified-preserved, 11 observed-partial, 9 altered, 1 confirmed-lost, 1 unresolved, and 5 measurement-error outcomes. Google Docs produced 11, 32, 11, 6, 5, and 0 respectively. A later full-corpus supplementary Microsoft Word evaluation adds 65 observations: 16 verified-preserved, 26 observed-partial, 16 altered, 2 confirmed-lost, no unresolved-equivalence, and 5 measurement-error outcomes. Six repeated variants were classification-stable and semantically stable 6/6 for each pipeline, although Word was byte-identical for 0/6. The 10 project-created realistic documents and 78 contracts per pipeline remain a separate applicability arm; Word adds 10 supplementary PDFs and 78 separately reported contract observations.

The evidence establishes condition-specific preservation behavior, not a global converter ranking or complete accessibility assessment. The strongest methodological lesson is that destination accessibility and preservation of known source semantics are different testing questions. A complete submission should therefore report the source contracts, destination evidence, uncertainty boundary, chronology of the two-pipeline primary design and the later Word supplement, and provenance together rather than infer preservation from tagged output alone.
